\def\SigleNumero{MAT-2901}
\def\NomCours{Mathématiques et technologie}
\def\Session{Hiver 2020}
\def\NomEvaluation{Travail pratique sommatif 1}
\def\Enseignant{Jérôme Soucy}
\def\EtudiantA{Lady Gaga}
\def\EtudiantB{Jean-Pierre Ferland}
\def\EtudiantC{Claude Dubois}
\def\EtudiantD{Marshall Bruce Mathers III}
\def\EtudiantE{Jennifer Lopez}
\documentclass{DevoirEtudiant}
	
\usepackage{commandesJS2017}	% Pour avoir accès aux commandes personnalisées

\begin{document}
	\begin{questions}
		\question Soit  $x\in\N$. Alors $x\in\Z, x\in\Q$ et $x\in\R$.
			\begin{parts}
				\part[12] Montrez que $x=2\Rightarrow x-2=0$.
					\begin{sol}
						Il suffit de soustraire $2$ de chaque membre de l'équation.
					\end{sol}
				\part[8] Trouvez les solutions dans $\R$ de l'équation $\sin (2x+\pi)=1$.
					\begin{sol}
						Nous avons que
						\begin{align}
							\sin (2x+\pi)=1 & \iff 2x+\pi=\frac{\pi}{2}+2k\pi\qquad(k\in\Z)\\
											& \iff 2x=\frac{\pi}{2}+2k\pi -\pi\\
											& \iff x=\frac{\pi}{4}+k\pi -\frac{\pi}{2}\\
											& \iff x=-\frac{\pi}{4}+k\pi.\\
						\end{align}
						Ainsi, l'ensemble solution est donné par $\left\{x\in\R:\exists\,k\in\Z \text{~vérifiant~} x=-\frac{\pi}{4}+k\pi\right\}$.
					\end{sol}
			\end{parts}
		\newpage
		\question[10] Voici comment inclure une image (matricielle) qui se trouve dans le même dossier que ce document.\\
		\includegraphics[width=6cm]{plan.PNG}
		\question[10] Voici une image créée avec Tikz, un outil pour créer des images vectorielles.\\
			\begin{tikzpicture}[scale=1]
			\draw[->] (-3,0) -- (3,0);
			\draw[->] (0,-3) -- (0,3);
			
			\draw node[right] (A) at (3,0) {\small $y$};
			\draw node[above] (B) at (0,3) {\small $z$};
			
			\foreach \x in {1,2}
			\draw[shift={(\x,0)}] (0pt,1pt) -- (0pt,-1pt) node[below] {\x};
			
			\foreach \y in {-1,1,2}
			\draw[shift={(0,\y)}] (-1pt,0pt) -- (1pt,0pt) node[left] {\y};
			\end{tikzpicture}\\
		Pour voir la différence entre les deux, zoomer.
	\end{questions}
\end{document} 